\section{Estrategia de testing}

\subsection{Descripcion}

Este documento define la estrategia general de pruebas del proyecto Quiniela. Su objetivo es asegurar que las reglas del dominio, los procesos criticos y los flujos visibles del producto se validen con el nivel de prueba adecuado segun su riesgo funcional y operativo.

\subsection{Alcance}

Aplica a pruebas unitarias, de integracion, end to end y validacion operativa. Cubre backend, frontend e integraciones en el nivel necesario para proteger reglas del negocio y confianza competitiva del sistema.

\subsection{Definiciones}

\begin{itemize}
  \item \textbf{Prueba unitaria}: validacion acotada de una regla o componente aislado.
  \item \textbf{Prueba de integracion}: validacion del comportamiento conjunto entre modulos, capas o contratos.
  \item \textbf{Prueba end to end}: validacion de un flujo funcional completo desde la perspectiva del usuario o del operador.
  \item \textbf{Prueba operativa}: validacion de procesos o contingencias relevantes para salida a operacion.
\end{itemize}

\subsection{Reglas}

\begin{itemize}
  \item El testing debe priorizar riesgo funcional antes que volumen de casos triviales.
  \item Las reglas que puedan producir disputa competitiva o error operativo deben tener cobertura prioritaria.
  \item Las pruebas deben distribuirse por nivel segun donde viva realmente el riesgo: dominio, contrato, flujo o integracion.
  \item Una validacion del frontend no reemplaza una prueba del enforcement real del backend.
\end{itemize}

\subsubsection{Prioridades de riesgo}

\begin{itemize}
  \item Prioridad alta: activacion de usuarios, lock de picks, scoring, resultados corregidos, recalculo, permisos administrativos y archivado.
  \item Prioridad media: visualizacion de leaderboard, feedback funcional, flujos de perfil y recuperacion de acceso.
  \item Prioridad baja: detalles de presentacion no normativos, siempre que no alteren acciones criticas.
\end{itemize}

\subsubsection{Distribucion recomendada por nivel}

\begin{longtable}{p{0.25\textwidth}p{0.25\textwidth}p{0.38\textwidth}}
\toprule
\textbf{Nivel} & \textbf{Enfoque} & \textbf{Que debe proteger} \\
\midrule
unitario & reglas acotadas y transformaciones puras & scoring, desempates, validaciones de rango, traduccion de estados, adaptacion de resultados \\
integracion & colaboracion entre modulos y contratos & activacion por comprobante, locks, sincronizacion de resultados, recalculo, permisos y respuestas de API \\
end to end & flujo visible completo & registro, activacion, captura de picks, bloqueo, leaderboard, admin y correccion manual \\
operativo & continuidad y contingencia & fallback manual, override, archivado y trazabilidad minima para operacion \\
\bottomrule
\end{longtable}

\subsubsection{Cobertura esperada por area}

\begin{itemize}
  \item Backend: enfasis en reglas del dominio, procesos criticos, validacion de frontera, idempotencia y contratos conceptuales.
  \item Frontend: enfasis en flujos visibles, estados de pantalla, validacion de experiencia y restricciones funcionales recibidas desde backend.
  \item Integraciones: enfasis en adaptacion de proveedor externo, prioridad de override manual y tolerancia a datos incompletos o tardios.
\end{itemize}

\subsubsection{Herramientas base aprobadas}

\begin{itemize}
  \item \texttt{Vitest} para pruebas unitarias y parte de integracion backend.
  \item \texttt{Playwright} para pruebas end to end del frontend y flujos principales.
  \item Validacion de contratos con \texttt{Zod} como apoyo al enforcement, sin sustituir pruebas funcionales.
\end{itemize}

\subsection{Edge Cases}

\begin{itemize}
  \item Un flujo puede requerir mas de un nivel de prueba para quedar realmente protegido; por ejemplo lock de picks debe existir en backend y en end to end.
  \item Un caso simple de UI puede volverse critico si comunica incorrectamente un estado de activacion, resultado corregido o restriccion por lock.
  \item Una integracion estable puede seguir necesitando pruebas de contingencia por su impacto operativo.
\end{itemize}

\subsection{Invariantes}

\begin{itemize}
  \item Ningun flujo critico se considera listo sin criterio de prueba verificable.
  \item Ninguna regla clave del dominio depende de una sola prueba de interfaz.
  \item Las pruebas deben poder explicar por que una regla critica esta protegida y en que nivel lo esta.
\end{itemize}

\subsection{Preguntas abiertas}

\begin{itemize}
  \item Confirmar el umbral minimo de cobertura formal que el proyecto exigira por repositorio y por nivel de prueba.
\end{itemize}

\subsection{Decisiones pendientes}

\begin{itemize}
  \item \texttt{Decision pendiente} \texttt{No bloqueante}: definir metas cuantitativas de cobertura por backend y frontend si el proyecto decide imponerlas.
  \item \texttt{Decision pendiente} \texttt{No bloqueante}: confirmar si las pruebas operativas de contingencia viviran dentro de suites automatizadas o como checklist manual controlado.
\end{itemize}
